\origpage{369}{N. Robertson, P.D. Seymour / Journal of Combinatorial Theory, Series B 90 (2004) 325--385}
\textbf{Subproof.} Since $(\Gamma \cap ins(D_v), \Gamma \cap out(D_v)) \in \mathscr{T} \setminus \mathscr{T}'$, it follows that $|E(F)| \le |E(D_v)|$, and by 13.5, $|E(D_v)| \le 4\rho + 4$.

\hspace*{12pt}(2) \textit{There is an edge of $\Gamma$ included in $ins(F) \cap ins(D_v)$.}

\textbf{Subproof.} $(\Gamma \cap out(F), \Gamma \cap ins(F)), (\Gamma \cap out(D_v), \Gamma \cap ins(D_v)) \in \mathscr{T}'$, and so not every edge belongs to $out(F) \cup out(D_v)$ by Theorem 2.3 of [5]. This proves (2).

\hspace*{12pt}(3) $ins(F) \cap ins(H) = U(F) \cap U(H)$.

\textbf{Subproof.} Otherwise either $|V(H)| = 1$ and $V(H) \subseteq ins(F) \setminus F$, or $H$ is a circuit and $ins(H) \subseteq ins(F)$, or $H$ is a circuit and $ins(F) \subseteq ins(H)$, by 8.6. The first and second contradict that $H$ is a $W$-hub (since $H \ne F$ by (2)), and the third contradicts (2). This proves (3).

\hspace*{12pt}(4) \textit{There is a closed disc $\Delta_0 \subseteq \hat{\Sigma}$ with $W \cup ins(F) \subseteq \Delta_0$, such that $\Delta_0$ includes $ins(C)$ for every circuit $C$ of $\Gamma^*$ with $|E(C)| \le 4\rho + 4$ and $U(C) \subseteq \Delta_0$.}

\textbf{Subproof.} If $ins(F) \subseteq W$ we may take $\Delta_0 = W$. Otherwise, since $U(H) \nsubseteq ins(F) \setminus U(F)$, it follows that $V(F) \cap bd(W) \ne \emptyset$. Choose $z^* \in V(F) \cap bd(W)$, and let $z \in A(\Gamma)$ with $z^* \in z$. By 7.1 (with $\kappa = 4\rho + 6$), since $\operatorname{rep}(\mathscr{T}) \ge 6\rho + 14$, there is a circuit of $sk(\Gamma)$ bounding an open disc $\Lambda \subseteq \hat{\Sigma}_\chi$ with $z \subseteq \Lambda$, such that

\hspace*{12pt}$\bullet$ $x \subseteq \Lambda$ for every $x \in A(\Gamma)$ with $d(z, x) \le 4\rho + 4$, where $d$ is the metric of $\mathscr{T}$, and

\hspace*{12pt}$\bullet$ $ins(C^*) \subseteq \bar{\Lambda}$ for every circuit $C^*$ of $\Gamma^*$ with $U(C^*) \subseteq \bar{\Lambda}$ and $|E(C^*)| < 4\rho + 4$.

Now if $T$ is a path of $\Gamma^*$ with $\le 4\rho + 4$ edges and with one end $z^*$, then $d(z, x) \le 4\rho + 4$ for every $x \in A(\Gamma)$ with $x \cap U(T) \ne \emptyset$, and so every such $x$ is a subset of $\Lambda$; and hence $U(T) \subseteq \Lambda$. But every edge of $H$ or of $\Gamma^* \cap bd(W)$ or of any spoke of $W$ belongs to some such path $T$, and hence $U(H) \subseteq \Lambda, bd(W) \subseteq \Lambda$ and $U(S) \subseteq \Lambda$ for every spoke $S$ of $W$. Let $\Delta_0 = \bar{\Lambda}$. Since $H$ and each $D_{v'}$ have at most $4\rho + 4$ edges, it follows that $ins(H), ins(D_{v'}) \subseteq \Delta_0$ for each $v'$, and so $W \subseteq \Delta_0$. This proves (4).

Since $ins(F), \Delta_v$ are both discs in $\Delta_0$ and their interiors intersect by (2), it follows from (4) that there is a circuit $C$ of $\Gamma^*$ with $C \subseteq F \cup D_v$, such that $U(C)$ bounds a closed disc in $\Delta_0$ including $ins(F) \cup \Delta_v$. \fix{To see this, consider the closed discs $\Delta \subseteq \Delta_0$ with $\Delta_v \subseteq \Delta$ that are bounded by $U(C')$ for some circuit $C'$ of $\Gamma^*$ with $C' \subseteq F \cup D_v$. There are only finitely many, since $F \cup D_v$ has only finitely many circuits and each of them bounds at most one disc in $\Delta_0$; and $\Delta_v$ is one of them. Choose such a disc $\Delta'$ maximal under inclusion, and let it be bounded by $U(C)$. It suffices to show that $ins(F) \subseteq \Delta'$. Each edge of $F$ not in $E(C)$ is disjoint from $U(C)$, which is the boundary of $\Delta'$, and so it is either included in $\Delta'$ or disjoint from $\Delta'$. If every edge of $F$ is included in $\Delta'$, then $U(F) \subseteq \Delta'$, and so the disc $ins(F)$ in $\Delta_0$ bounded by $U(F)$ is included in $\Delta'$. If no edge of $F$ is included in $\Delta'$, then $U(F)$ is disjoint from the interior of $\Delta'$; since this interior is connected and meets the interior of $ins(F)$ (because it includes the interior of $\Delta_v$, and by (2)), it is included in the interior of $ins(F)$, and so $\Delta' \subseteq ins(F)$. But then $ins(F)$ is one of our discs, and so $ins(F) = \Delta'$ by the maximality of $\Delta'$. Finally, suppose that some edge $e_0$ of $F$ is disjoint from $\Delta'$ and some edge $e_1$ of $F$ is included in $\Delta'$. Then both ends of $e_1$ belong to $\Delta'$, and so there is a path $P \subseteq F$ with $e_0 \in E(P)$, with distinct ends $u, v \in \Delta'$ and with no other vertex in $\Delta'$. Each edge of $P$ is disjoint from $\Delta'$, since it is $e_0$ or has an end not in $\Delta'$; hence $U(P) \cap \Delta' = \{u, v\}$, and $u, v$ belong to the boundary $U(C)$ of $\Delta'$, so $u, v \in V(C)$. Since $P$ is a path in $\Delta_0$ between two vertices of $C$ which meets $\Delta'$ only in its ends, one of the two circuits of $C \cup P$ containing $P$ bounds a disc in $\Delta_0$ which is the union of $\Delta'$ and the disc in $\Delta_0$ bounded by the other one, and so properly includes $\Delta'$. This contradicts the maximality of $\Delta'$, and so proves the existence of $C$.}\fixnote{fix:13.6-C-proof-added}{The paper does not say how the circuit $C$ is found. The argument above is the one used in the formalization.}

\hspace*{12pt}(5) $|E(C)| \le |E(D_v)|$, \textit{and if equality holds then} $\lambda(C) \le \lambda(D_v)$.

\textbf{Subproof.} Let $(A, B) = (\Gamma \cap out(F), \Gamma \cap ins(F))$, and $(A', B') = (\Gamma \cap out(D_v), \Gamma \cap ins(D_v))$. Since $(A, B), (A', B') \in \mathscr{T}'$ it follows from the second tangle axiom that $(B \cap B', A \cup A') \notin \mathscr{T}'$. But the latter has order at most the sum of the orders of $(A, B)$ and $(A', B')$, and hence at most $4\rho + 4$ by (1), and so $(B \cap B', A \cup A') \in \mathscr{T}$ since $(B, A) \in \mathscr{T}$. Thus $(B \cap B', A \cup A') \in \mathscr{T} \setminus \mathscr{T}'$, and it follows that $(A \cup A', B \cap B')$ does not have smaller $\lambda$-order than $(A, B)$. By 8.3, $(A \cap A', B \cup B')$ has $\lambda$-order at most that of $(A', B')$. In particular, since every edge of $C$ is split by $(A \cap A', B \cup B')$ and only edges
